% ==============================================================================
% T0 Theory / FFGFT: Dokument 284 (Kapitelinhalt, DE)
% Titel: Eine Formel-genaue Lesung des Spiral-Zeit-Neurodynamik-Modells
% Autor: Johann Pascher
% Datum: 19. Juni 2026
% Zweck: Reproduzierbarer, formel-genauer Test von Krueger/Feeney/Wende
%        (Medinformatics 2026); das grundlegende Zeitachsen-Problem; die FFGFT-
%        Lesung Kompaktifizierung / Umwandlung-als-Messung.
% Referenzen: Dok. 190, 210, 249, 263, 270, 282, 283.
% Skripte: test_algebra.py, test_deltaphi.py, test_dynamics.py,
%          test_compactification.py, make_figures.py
% ==============================================================================

\section*{Kurzfassung}
Dieses Dokument ist eine reproduzierbare, formel-genaue Lesung von \emph{Information-Theoretic
Modeling of Neural Coherence via Triadic Spiral-Time Dynamics: A Framework for
Neurodynamic Collapse} (Kr\"uger, Feeney \& Wende, \emph{Medinformatics} 2026,
DOI 10.47852/bonviewMEDIN62029043). Jeder Befund wird von einem kurzen, lauffähigen
Skript erzeugt. Die Lesung trennt, was trägt (ein interpretierbares
Phasen-Abweichungsmerkmal; echte Oktonion-Nichtassoziativität), von dem, was falsch
etikettiert ist (die Tripel-Algebra ist quaternionisch, nicht bikomplex), intern
widersprüchlich (zwei unvereinbare Definitionen von $\Delta\Phi$), nicht hergeleitet
(die $0{,}40$-Schwelle; die behaupteten Spiralen und Grenzzyklen) oder aus dem Papier
nicht reproduzierbar (das berichtete $r\approx0{,}96$). Darunter liegt \emph{ein}
grundlegendes Problem: der Operator $\psi(t)=t+i\varphi(t)+j\chi(t)$ ist auf der
physikalischen, linearen Zeit als Basisachse aufgebaut -- der dekompaktifizierten,
nach-der-Messung-Koordinate -- sodass sein Gedächtnis per Konstruktion generisch ist.
Die FFGFT-Lesung macht das präzise: das System ist kompakt und trägt Gedächtnis. Nach
Dok.~270 ist die Dekompaktifizierung in die SI-Zeit ein mode-erhaltendes Embedding
(verlustfrei), und FFGFT hat sehr wohl framework-spezifische SI-Observable
(Massen/Koide, CHSH$\sim\xi$, $D_f$, $L_0$; Dok.~270/249); was die Wiederkehr entfernt,
ist nicht die SI-Umwandlung, sondern das endliche, nichtperiodische Beobachtungsfenster
plus Dekohärenz. Die scharfe, enge Folge: die Kohärenz-Sektor-Recurrence ist auf einem
warmen, endlich-gefensterten SI-Zeit-Signal wie EEG nicht \emph{auflösbar}. Dies
ist eine Gutachter-Notiz an die Autoren; der Maßstab ist die Falsifikations-Disziplin,
die die HLV/HICE-Seite bereits anwendet.

% ==============================================================================
\section{Anlass und Methode}
% ==============================================================================
Das Medinformatics-Papier wendet den HLV-Spiral-Zeit-Operator auf neuronale Kohärenz an
und schlägt das Information--Kohärenz--Energie-Prinzip (ICE) sowie eine
$\Delta\Phi$-gestützte ``neurodynamische Kollaps''-Grenze vor. Statt aus dem Eindruck zu
argumentieren, wird hier jede mathematische Behauptung rechnerisch geprüft. Die Skripte
sind \texttt{test\_algebra.py}, \texttt{test\_deltaphi.py}, \texttt{test\_dynamics.py},
\texttt{test\_compactification.py} und \texttt{make\_figures.py}; durchgängig feste Seeds.
Die Befunde teilen sich in: \emph{trägt}, \emph{falsch etikettiert}, \emph{intern
widersprüchlich}, \emph{nicht hergeleitet}, \emph{aus dem Papier nicht reproduzierbar}
und \emph{spekulative Prämisse}. Ein Problem ist grundlegend -- die Zeitachse, auf der der
Operator aufgebaut ist -- und wird zuerst behandelt, weil die Befunde zu $\Delta\Phi$ und
zur Dynamik seine Symptome sind.

\begin{center}\small
\renewcommand{\arraystretch}{1.25}
\begin{tabular}{@{}c p{4.1cm} p{1.7cm} p{5.0cm}@{}}
\toprule
Nr. & Behauptung & Abschnitt & Befund\\
\midrule
0 & Basis-Zeitachse $\psi(t)=t+i\varphi(t)+j\chi(t)$ auf physikalischer Linearzeit & 2.1, 2.2.1 & Grundlegend -- Operator auf der nach-der-Messung-Achse (ordnet 3 \& 5 unter)\\
1 & Tripel-Zeit ``reduzierte bikomplexe/oktonionische Ebene'' & 2.1 & Falsch etikettiert -- es sind die Quaternionen\\
2 & Oktonionischer Operator; Nichtassoz.\ ``erklärt Psychologie'' & 2.10 & Mathematik trägt, Behauptung überdehnt\\
3 & $\Delta\Phi$-Instabilitätsfunktional & 2.6 vs 12/15 & Intern widersprüchlich (zwei Definitionen)\\
4 & kritisch $\Delta\Phi_c\approx0{,}40$, ``kein freier Parameter'' & 2.5, 2.6.2 & Nicht hergeleitet (freier Parameter)\\
5 & gesunde Spiralen / Trauma-Grenzzyklen & 2.2.3, 2.9 & Nicht aus den Gleichungen hergeleitet\\
6 & $r\approx0{,}96$, $\Delta\Phi_c\approx0{,}40$ ``über Probanden'' & 3.4 & Aus dem Papier nicht reproduzierbar\\
7 & Axon-Biphotonen-Kopplung $\Gamma_{\text{axon}}$ & 2.7--2.8 & Spekulative Prämisse als etabliert\\
\bottomrule
\end{tabular}
\end{center}

% ==============================================================================
\section{Die Basis-Zeitachse (grundlegend)}
% ==============================================================================
Dies ist die folgenreichste Wahl im Papier und jene, unter der die Befunde 3 und 5
nachgeordnete Symptome sind. Der Spiral-Zeit-Operator
\begin{equation}
\psi(t) = t + i\,\varphi(t) + j\,\chi(t)
\end{equation}
nimmt die physikalische, lineare Zeit $t$ als Skalarbasis und schreibt $\varphi,\chi$ als
Funktionen ebendieser $t$. Das Papier benennt $t$ ausdrücklich als physikalisch: $U_1$
ist als ``lineare physikalische Zeit'', äußere Chronologie definiert (\S2.2.1). Das legt
die Darstellung fest, in der jede spätere Gedächtnis- und Holonomie-Behauptung getroffen
wird.

\paragraph{Einheiten-unabhängig.} Damit $t + i\varphi + j\chi$ \emph{ein} algebraisches
Objekt ist, müssen die drei Komponenten dieselbe Art von Größe sein. Das Papier behauptet
Dimensionslosigkeit (\S2.1.2), nennt $t$ aber zugleich physikalische Zeit und Chronologie
(\S2.2.1); $\varphi$ und $\chi$ sind interne Phasen-/Identitätskoordinaten. Eine
physikalische Zeit zu internen Phasen in einer quaternionischen Summe zu addieren ist
dimensionell inkohärent, solange nicht alle drei zuvor auf ein gemeinsames dimensionsloses
Maß reduziert werden. (Da $i,j$ antikommutieren, ist das Objekt quaternionisch -- $t$ ist
der Realteil eines Quaternions, dessen Kreuzterm $ij=k$ nie mitgeführt wird; siehe
\S\ref{sec:alg}.)

\paragraph{Tieferes.} Ein framework-spezifisches Gedächtnis lebt auf dem kompakten
internen Zeitkreis $S^1$: ein diskretes Spektrum, eine exakte Wiederkehr (in FFGFT der
Back-flow $5{,}125$ im Kohärenz-Sektor, Dok.~282). $\varphi(t),\chi(t)$ als einwertige
Funktionen der physikalischen linearen $t$ zu schreiben, legt $S^1$ auf $\mathbb{R}$ flach
-- es führt die Projektion $S^1\!\to\!\mathbb{R}$ \emph{vor jeder Dynamik} aus, und eine
einwertige Funktion über linearer Zeit hat keine kompakte Periodizität, sodass die
Wiederkehr auf dieser Achse gar nicht darstellbar ist. Nach dem Prinzip
Umwandlung-als-Messung (\S\ref{sec:ffgft}) ist die Festlegung auf physikalische Zeit selbst
die störende, informationsverwerfende Projektion -- kein Apparat nötig. Die Schreibweise
hat diesen Schritt stillschweigend vollzogen und verbirgt die eine Operation -- die
Umwandlung/Messung von der kompakten internen Zeit auf die physikalische SI-Zeit --, die
entscheidet, ob eine Gedächtnis-Behauptung framework-spezifisch oder generisch ist.

\paragraph{Folge.} Der Operator ist per Konstruktion nach-der-Messung: was immer er auf
physikalischer (EEG-)Zeit an Gedächtnis trägt, ist generisch. Das ist reproduzierbar
(\texttt{test\_compactification.py}): das Rückkehr-Signal des diskreten Spektrums hat auf
dem kompakten Horizont den Back-flow $11{,}5$, den ein endliches physikalisches
Auslesefenster auf $0{,}000$ treibt, während die sechs Linien nie auflösbar werden
($2\pi/T_{\text{Fenster}}\gg$ Linienabstand $0{,}091$) -- die Auslesung ist eine
nicht-invertierbare Projektion. Dieselbe Wurzel erklärt die beiden nachgeordneten Befunde:
eine Berry-Phasen-Holonomie (Gl.~12/15, \S\ref{sec:dphi}) braucht eine geschlossene
Schleife -- einen kompakten Zyklus --, und ein Grenzzyklus (\S\ref{sec:dyn}) braucht
nicht-gradiente geschlossene Orbits; beides ist auf einer linearen physikalischen
Zeitachse nicht tragfähig.

\begin{figure}[h]\centering
\includegraphics[width=0.95\linewidth]{fig_compactification.png}
\caption{Auslesung als nicht-invertierbare Projektion. Links: in der kompakten Darstellung
erzeugt das diskrete Spektrum Wiederkehren (Back-flow $\gg0$); das rote Band markiert ein
endliches physikalisches Auslesefenster. Rechts: im Fenster zerfällt das Observable ohne
Wiederkehr, und die diskreten Linien sind unaufgelöst.}
\end{figure}

\paragraph{Konstruktiver Ausweg.} Den Operator auf die kompakte, dimensionslose interne
Zeit (zyklisch) stellen, auf der das diskrete Spektrum und die Wiederkehr leben, und die
physikalische SI-Zeit nur über eine ausdrücklich deklarierte Umwandlung eintreten lassen
-- und genau diese projiziert den framework-spezifischen Anteil heraus. Entweder bleibt
der Operator intern/kompakt (keine direkte EEG-Aussage), oder er ist physikalisch/SI
(generisches Gedächtnis). $\psi(t)=t+i\varphi(t)+j\chi(t)$ versucht beides. Die minimale
Änderung ist, zwei verschiedene Zeiten zu deklarieren und die Umwandlung zwischen ihnen
ausdrücklich als die Messung zu schreiben, die sie ist.

% ==============================================================================
\section{Tripel-Zeit-Algebra -- falsch etikettiert}\label{sec:alg}
% ==============================================================================
Die angegebenen Relationen $i^2=j^2=-1,\ ij=-ji$ sind genau die imaginären Quaternionen:
mit $k:=ij$ ist $k^2=-1$, und $\{1,i,j,k\}$ schließt zu $\mathbb{H}$, das assoziativ und
nichtkommutativ ist (\texttt{test\_algebra.py}, A). Eine \emph{bikomplexe} Ebene verlangt,
dass die beiden imaginären Einheiten kommutieren ($ij=ji$); hier gilt $ij=-ji$, also ist
die Algebra \emph{nicht} bikomplex. Ein Drei-Term-Operator $t+i\varphi+j\chi$ lebt in einer
Quaternionen-Unteralgebra; Oktonionen sind dafür nicht nötig. Der korrekte Name ist
quaternionisch.

% ==============================================================================
\section{Oktonionische Struktur -- Mathematik trägt, Deutung überdehnt}
% ==============================================================================
Die Oktonion-Nichtassoziativität ist real: der Assoziator
$(e_1e_2)e_4-e_1(e_2e_4)=2e_7\neq0$ (Cayley--Dickson), während Assoziatoren auf einem
quaternionischen Tripel verschwinden (\texttt{test\_algebra.py}, B). Das ist korrekte
Mathematik. Der weitere Schritt -- dass dies ``empirische Psychologie direkt abbildet'' und
``die erste physikbasierte Erklärung psychologischer Nichtlinearität'' sei -- ist eine
interpretierende Analogie, keine Herleitung: keine psychologische Größe wird aus der
Algebra berechnet oder gegen sie getestet, und die acht Basisrichtungen werden per
Festsetzung mit Bedeutungen belegt.

% ==============================================================================
\section{$\Delta\Phi$ -- zwei unvereinbare Definitionen unter einem Symbol}\label{sec:dphi}
% ==============================================================================
\begin{itemize}
\item \textbf{Definition 1 (\S2.6):} $\Delta\Phi=\alpha|\Delta S|+\beta|\Delta I|+\gamma|\Delta C|$,
konvexe Gewichte, jede Komponente in $[0,1]$ -- ein einheitenloser gewichteter
$L^1$-Index in $[0,1]$. Er ist in jeder Komponente linear, hat also keine Krümmung und
keine intrinsische Bifurkation; die Niveaumenge $\Delta\Phi=c$ ist für jedes $c$ eine
flache Polytop-Facette.
\item \textbf{Definition 2 (Gl.~12, 15):} $\Delta\Phi=\oint(\nabla\times A)\cdot da
\approx \mathrm{Im}\oint\langle\psi|\nabla\theta|\psi\rangle\,d\theta$ -- eine
Berry-Phasen-Holonomie, ein Winkel im Bogenmaß über $0\ldots2\pi$, ohne Grund, in $[0,1]$
oder nahe $0{,}40$ zu liegen (eine Schleife mit $|\gamma|=0{,}40$\,rad ist ein winziger
$\sim$\,29$^\circ$-Kegel).
\end{itemize}
Ein normierter Abweichungsindex und eine Holonomie im Bogenmaß sind verschiedene
mathematische Objekte. Beide ``$\Delta\Phi$'' zu nennen und beiden denselben kritischen
Wert $0{,}40$ zuzuweisen, ist eine Äquivokation (\texttt{test\_deltaphi.py}, C/D).

\begin{figure}[h]\centering
\includegraphics[width=0.72\linewidth]{fig_deltaphi.png}
\caption{Die zwei $\Delta\Phi$-Definitionen sind unvereinbare Objekte: eine Holonomie im
Bogenmaß ($0\ldots2\pi$) gegenüber einem einheitenlosen Index ($0\ldots1$).}
\end{figure}

% ==============================================================================
\section{Die $0{,}40$-Schwelle -- nicht hergeleitet}
% ==============================================================================
\S2.5 sagt, der Wert sei ``kein freier Tuning-Parameter'' und die Operatorstruktur
``erzwinge'' eine begrenzte Instabilitätsfläche. Doch das $L^1$-Funktional hat keinen
intrinsischen kritischen Punkt (es ist linear); die einzig verfügbare strukturelle Aussage
ist die triviale Schranke $0\le\Delta\Phi\le1$ und die symmetrische Referenz
$\Delta\Phi=d$ (also $1/3$ bei $d=1/3$). \S2.6.2 selbst schreibt die Differenz $0{,}40$ vs
$0{,}333$ ``datensatzspezifischer Gewichtungs-Asymmetrie $\ldots$ und Normierung'' zu, und
bei gleicher Abweichung $d=0{,}40$ liefert jede konvexe Gewichtung $0{,}40$
(\texttt{test\_deltaphi.py}, C/E). Also ist $0{,}40$ ein empirischer/Konturwert, keine
strukturell erzwungene Bifurkation.

% ==============================================================================
\section{Spiralen und Trauma-Grenzzyklen -- nicht hergeleitet}\label{sec:dyn}
% ==============================================================================
Der einzige mit expliziten Potenzialen angegebene dynamische Mechanismus (\S2.4) ist der
Gradientenfluss-Teil der Mastergleichung, $\partial_t\Psi=-\nabla V_{\text{DLHR}}(\Psi)$.
Für jeden Gradientenfluss gilt
\begin{equation}
\frac{dV}{dt}=\nabla V\cdot\frac{d\Psi}{dt}=-|\nabla V|^2\le0,
\end{equation}
also ist $V$ eine strikte Ljapunow-Funktion, und das System kann keine periodischen Orbits,
Grenzzyklen oder anhaltenden Spiralen haben -- es relaxiert in Fixpunkte (Theorem).
Numerisch (\texttt{test\_dynamics.py}, F): über 40 Trajektorien ist $V(t)$ monoton fallend
(0 Verletzungen) und 0 geschlossene Orbits werden gefunden. Die Phänomenologie ``gesunde
Kognition $\to$ stabile Einwärtsspiralen'', ``Trauma-Schleifen $\to$ Grenzzyklen'' folgt
daher nicht aus der spezifizierten Dynamik; sie bräuchte die Advektionsterme
$\lambda_\varphi\partial_\varphi,\ \lambda_\chi\partial_\chi$ oder $W_{\text{oct}}$, die nie
konkret angegeben werden. (Der kubische $M^3$-Term macht $V$ zudem nach unten
unbeschränkt.)

\begin{figure}[h]\centering
\includegraphics[width=0.95\linewidth]{fig_dynamics.png}
\caption{Gradientenfluss mit den DLHR-Potenzialen des Papiers: Trajektorien relaxieren in
Fixpunkte (keine geschlossenen Orbits, keine anhaltende Spirale), und $V(t)$ fällt monoton,
sodass Grenzzyklen unmöglich sind.}
\end{figure}

% ==============================================================================
\section{Empirische Behauptungen -- aus dem Papier nicht reproduzierbar}
% ==============================================================================
Das berichtete $r\approx0{,}96$ und $\Delta\Phi_c\approx0{,}40$ ``über Probanden'' (\S3.4)
lassen sich aus dem Papier nicht nachrechnen, das keine Datentabelle gibt. Die Reproduktion
braucht den externen Datensatz (Zenodo 10.5281/zenodo.17990696) und das Validierungs-Notebook.
Ein fairer Test bräuchte den Datensatz plus ein vorregistriertes Protokoll mit den Nullmodellen,
die die Autoren auf der HICE-Seite bereits verwenden (Zeit-Shuffle, present-state-matched,
phasen-randomisiert).

% ==============================================================================
\section{Axon-Biphotonen-Verschränkung -- spekulative Prämisse als etabliert}
% ==============================================================================
$\Gamma_{\text{axon}}=g\,e^{i\varphi}E_{12}+g'e^{j\chi}E_{23}$ (\S2.8) ruht auf
``experimentellen Befunden'' von Biphotonen-Verschränkung und Langstrecken-Kohärenz in
myelinisierten Axonen, dargestellt als gesicherte Forschung. Verschränkung in warmen,
nassen, myelinisierten Axonen ist keine etablierte Physik. Das ist kein Formelfehler, aber
eine tragende empirische Prämisse, die als Tatsache behauptet wird und als Hypothese zu
kennzeichnen ist.

% ==============================================================================
\section{Ein verteidigbarer Kern}
% ==============================================================================
Ohne den Überbau bleibt ein solides, bescheidenes Objekt:
$\Delta\Phi=|\omega(t)-\omega_{\text{memory}}(t)|$ -- die Abweichung der momentanen
Phasengeschwindigkeit von einem kurzen Phasengedächtnis -- ist ein interpretierbares,
leichtes Zeitreihen-Instabilitätsmerkmal, und als Signalanalyse-Operator (die Rahmung des
begleitenden Atem-Papiers) legitim. Die Probleme beschränken sich auf den
HLV-/Oktonionen-/Bewusstseins-Überbau, die $\Delta\Phi$-Umdefinition sowie die nicht
hergeleitete Schwelle und Dynamik.

% ==============================================================================
\section{Eine FFGFT-Lesung: Umwandlung als Messung}\label{sec:ffgft}
% ==============================================================================
Zeitliche Rekursion trägt Gedächtnis, und das bildet auf FFGFTs eigene Struktur ab (der
Rekursionsoperator $\mathcal{R}$; der Gedächtniskern als Fourier-Transformierte der
Spektraldichte des Verbindungs-Laplace-Operators, Dok.~282). Das Folgende ist als
konstruktive Brücke und ausdrückliche Vermutung gemeint, nicht als Resultat.

FFGFTs Gedächtnis hat zwei Gesichter. Auf dem überdämpften / Relaxations-Observablen ist der
Kern eine positive Multiexponentielle, und ein generischer Biexponential reproduziert ihn
(Residuum $\sim10^{-3}$): nicht-identifizierbar, Stufe-0. Auf dem Kohärenz-Observablen
$c(t)=e^{-\Gamma(t)}$ -- seinem Heimat-Sektor -- erzeugt das diskrete Spektrum eine echte
Wiederkehr: Back-flow $5{,}125$, den generische zerfallende Kerne nicht reproduzieren
(Back-flow $0$). Das Gedächtnis ist also diskriminierend, aber nur im Kohärenz-Sektor und
nur, wenn \emph{beides} zugleich gilt: ein diskretes Spektrum \emph{und} erhaltene Kohärenz
(kleines $\Gamma$).

FFGFT ist fundamental kompakt ($T^4$, diskretes Spektrum), und in dieser dimensionslosen
Darstellung hat das System immer Gedächtnis. Es ist wichtig, den Verlust präzise zu
verorten, denn die Umwandlung in SI ist \emph{nicht} die Stelle, wo er geschieht. Nach
Dok.~270 ist die Typ-I-Dualitäts-Dekompaktifizierung $S^1_m\!\to\!\mathbb{R}_t$ (geregelt
durch $\tilde{T}\cdot m=1$) ein \emph{mode-erhaltendes Embedding} -- verlustfrei; die
Periodizität ist bloß nicht mehr direkt sichtbar und überlebt als diskretes Spektrum und
als physikalische Observable (die Revival-Zeit $T_\text{rev}$, die CHSH-Abweichung
$\sim\xi$). Eine reine Einheiten-Skalierung (mit $\hbar,c$) ist ebenso invertierbar und
verliert nichts. Die wirklich verlustbehafteten Operationen sind drei und verschieden: die
\emph{Rückrichtung} der Kompaktifizierung $\mathbb{R}\!\to\!S^1$ (nichtperiodische Anteile
steigen nicht herab), die räumlichen Projektionen $D3\!\to\!D0$ (Dok.~270) und -- für jede
Auslesung -- die Beschränkung auf ein \emph{endliches, nichtperiodisches Beobachtungsfenster
plus Dekohärenz}. Es ist dieser letzte Schritt, nicht die Einheiten-Umwandlung, der das
lokale Kurzzeitverhalten behält und die globale Wiederkehr verwirft -- genau wie eine
Messung eine Basis wählt und die Kohärenzen fallen lässt (Test~G). Eine physikalische
Messung ist dieselbe Fenster-und-Dekohärenz-Projektion an ihrem Starkkopplungs-Ende. Das ist
verträglich mit dem, was der Korpus bereits führt: die Vakuumenergie ist kompaktifiziert
vorhanden, dekompaktifiziert sichtbar, mit $c,\hbar c$ als SI-Umrechnungsfaktoren (Dok.~263)
-- die dekompaktifizierte Sicht ist nicht weniger wahr, nur anders aufgelöst.

Die Lesung stützt also den bescheidenen Kern (Rekursion trägt Gedächtnis) und begrenzt die
Behauptung präzise. Die framework-spezifische Signatur überlebt sehr wohl in die SI-Welt --
in den richtigen Kanälen: Mass-reading, CHSH$\sim\xi$, $D_f$, $L_0$ (Dok.~270/249). Was
\emph{nicht} überlebt, ist die Kohärenz-Sektor-Recurrence, gelesen auf einem warmen,
endlich-gefensterten SI-Zeit-Observablen: dort haben Fenster und Dekohärenz die globale
Wiederkehr bereits verworfen, sodass gerade das Observable, das die klinischen Abschnitte
verwenden -- SI-Zeit-EEG --, in genau diesem Sinn nach-der-Messung ist und das Regime
darstellt, in dem der FFGFT-Harness (Dok.~282/283) das Gedächtnis als generisch zeigt, nicht
als Fingerabdruck. Auch deshalb ist die Wahl der Basis-Zeitachse (\S3) entscheidend: der
Operator wird von Anfang an auf dieser endlichen, nach-der-Messung-Achse geschrieben.

% ==============================================================================
\section{Reproduzierbarkeit}
% ==============================================================================
\begin{itemize}
\item \texttt{test\_algebra.py} -- Quaternionen-Abschluss + Bikomplex-Widerlegung (A); Oktonion-Assoziator (B).
\item \texttt{test\_deltaphi.py} -- $L^1$-Schranke/Linearität (C); Berry-Phasen-Bereich (D); $0{,}40$ als freier Knopf (E).
\item \texttt{test\_dynamics.py} -- Gradientenfluss-Ljapunow / Test ohne Grenzzyklen (F).
\item \texttt{test\_compactification.py} -- Auslesung als nicht-invertierbare Projektion: kompakter Back-flow vs.\ Fenster (G).
\item \texttt{make\_figures.py} -- die oben verwendeten Abbildungen.
\end{itemize}
Die Abbildungen \texttt{fig\_deltaphi.png}, \texttt{fig\_dynamics.png} und
\texttt{fig\_compactification.png} werden von diesen Skripten erzeugt; alle mit festen Seeds.
